%% kapa: exam=1 round=25 year=2014 period=2 form=A subject=accounting subject_ko=회계학 seq=30 number=70
%% source: Q-Net 감정평가사 기출문제 (cst00309) artlSeq=5101207 file=extracted/1008/5101207/2014년 제25회 감정평가사 1차시험 기출문제/2교시(A형).hwp
%% conversion: kordoc 4.24.51 -> tools/gen1.py (auto); review=proofread-sonnet; flags=none
\begin{question}{70}
(주)서울은 영업 첫해인 20×1년의 법인세비용차감전순이익은 ￦800,000이고 과세소득은 ￦1,200,000이며, 이 차이는 일시적 차이로서 향후 2년간 매년 ￦200,000씩 소멸될 것이다. 20×1년과 20×2년의 법인세율은 40\%이고 20×1년에 개정된 세법에 따라 20×3년부터 적용될 법인세율은 35\%이다. (주)서울이 이 차이에 관하여 20×1년 말 재무상태표 상에 기록하여야 하는 이연법인세자산 또는 이연법인세부채의 금액은? (단, 이연법인세자산 또는 이연법인세부채는 각각 자산과 부채의 인식요건을 충족한다.)
\begin{choices}
  \item 이연법인세자산 ￦140,000
  \item 이연법인세부채 ￦140,000
  \item 이연법인세자산 ￦150,000
  \item 이연법인세부채 ￦150,000
  \item 이연법인세자산 ￦160,000
\end{choices}
\end{question}
\answer{3}
